% trans_minus29_full.tex
% The Trans -2/9 Identity for Degree-(2,1) Polynomial Continued Fractions:
% Empirical Evidence, an Algebraic Signature, and the Limits of Standard
% Asymptotic Methods.
%
% Companion to t2b_conjecture_note.tex / t2b_monat_submission_R1.tex
% (Zenodo DOI 10.5281/zenodo.19801038).
%
\documentclass[11pt]{amsart}

\usepackage[T1]{fontenc}
\usepackage{lmodern}
\usepackage{microtype}
\usepackage{amsmath,amssymb,amsthm}
\usepackage{mathtools}
\usepackage{booktabs}
\usepackage{enumitem}
\usepackage{hyperref}
\hypersetup{colorlinks=true, linkcolor=blue, citecolor=blue, urlcolor=blue}

\newtheorem{theorem}{Theorem}[section]
\newtheorem{proposition}[theorem]{Proposition}
\newtheorem{lemma}[theorem]{Lemma}
\newtheorem{corollary}[theorem]{Corollary}
\newtheorem{conjecture}[theorem]{Conjecture}
\newtheorem{observation}[theorem]{Observation}
\theoremstyle{definition}
\newtheorem{definition}[theorem]{Definition}
\theoremstyle{remark}
\newtheorem{remark}[theorem]{Remark}
\newtheorem{problem}{Problem}

\newcommand{\FF}{\mathcal{F}}
\newcommand{\KK}{\mathcal{K}}
\newcommand{\Rat}{\mathrm{Rat}}
\newcommand{\Log}{\mathrm{Log}}
\newcommand{\Alg}{\mathrm{Alg}}
\newcommand{\Trans}{\mathrm{Trans}}
\newcommand{\Des}{\mathrm{Des}}
\newcommand{\dps}{\mathrm{dps}}
\newcommand{\Z}{\mathbb{Z}}
\newcommand{\Q}{\mathbb{Q}}
\newcommand{\R}{\mathbb{R}}
\newcommand{\C}{\mathbb{C}}

\emergencystretch=1.5em

\title[The Trans $-2/9$ Identity]{The Trans $-2/9$ Identity for
Degree-$(2,1)$ Polynomial Continued Fractions:\\
Empirical Evidence, an Algebraic Signature, and the Limits of
Standard Asymptotic Methods}

\author{Papanokechi}
\address{Independent researcher, Yokohama, Japan}
\email{shkubo@outlook.jp}

\subjclass[2020]{11A55, 11J70, 11J81, 11Y65, 39A06, 40A15}

\keywords{Polynomial continued fractions, arithmetic stratification,
indicial equation, Birkhoff asymptotics, transcendental numbers,
PSLQ, integer-relation algorithms, experimental mathematics}

\begin{document}

\begin{abstract}
For an integer polynomial continued fraction (PCF) of partial-numerator
degree~$2$ and partial-denominator degree~$1$, write
$a(n) = a_2 n^2 + a_1 n + a_0$ and $b(n) = b_1 n + b_0$ with
$a_2,b_1\ne 0$, and set the Worpitzky parameter $r := a_2/b_1^2$.
A computational survey across approximately $585{,}000$ integer
families finds zero counterexamples to the empirical statement that
every transcendental-stratum (Trans) limit lying in the
negative-Worpitzky interior arises at $r = -2/9$, equivalently at
indicial exponents $\{1/3,2/3\}$ for the standard
plus-form recurrence; this is the
\emph{Transcendental Ratio Identity} announced in the companion
note~\cite{TRINote}.  The present paper extends that announcement
along two axes.

First, we establish the \emph{algebraic signature} of the locus
$r=-2/9$: under the alternative minus-form recurrence
$y_{n+1} = b(n)y_n - a(n)y_{n-1}$, the characteristic roots are
$\mu_\pm = b_1(3\pm\sqrt{17})/6$, and we prove that, for $a_2<0$,
$b_1\ne 0$, $r=-2/9$ is equivalent to the modulus identity
$|\mu_+/\mu_-| = (13+3\sqrt{17})/4$ (PSLQ-certified at $\dps=150$
with integer relation $[13,3,-4]$ on $\{1,\sqrt{17},\rho\}$).  This
is the first explicit identification of a quadratic irrational
($\sqrt{17}$) as an algebraic invariant of the Trans property.

Second, we record three negative results that bound the
analytic landscape: (i) under the minus-form recurrence the
indicial pair $\{1/3,2/3\}$ corresponds to a compatibility ideal
intersecting $r=-2/9$ only at irrational $b_0/b_1=(27\pm\sqrt{17})/18$,
so no integer family realises both simultaneously; (ii) the leading
two-term Birkhoff balance produces no rational-coefficient invariant
that distinguishes the Trans locus beyond the (Vieta-tautological)
ratio identity; (iii) the sub-leading $[n^{-1}]$ Birkhoff balance is
regular on the Trans locus, with the candidate resonance coefficient
factoring as $(a_2-b_1\mu)(4a_2-b_1^2)$ and vanishing only on the
algebraic-discriminant-zero locus, never on Trans.

The combined picture is that the Trans property of degree-$(2,1)$ PCFs
is empirically robust, algebraically fingerprinted by~$\sqrt{17}$, and
not accessible to leading-order or single-step sub-leading Birkhoff
methods.  We pose the conjecture as a problem in the Galois theory of
linear difference equations and in isomonodromic deformation theory.
\end{abstract}

\maketitle

%% ==================================================================
\section{Introduction}\label{sec:intro}
%% ==================================================================

A \emph{polynomial continued fraction (PCF) of degree $(d,e)$ with
integer coefficients} is an expression
\begin{equation}\label{eq:pcf}
  K(a,b) \;=\; b_0 \,+\, \cfrac{a(1)}{b(1) \,+\, \cfrac{a(2)}{b(2)
  \,+\, \cdots}}\,,
\end{equation}
where $a(n) = a_d n^d + \cdots + a_0$ and $b(n) = b_e n^e + \cdots +
b_0$ are integer polynomials with $a_d \ne 0$ and $b_e \ne 0$.  The
case $(d,e) = (2,1)$ is the smallest non-trivial case in which
classical limits such as $\pi$, Brouncker's $4/\pi$, and Apéry-like
series all appear.  Writing $a(n) = a_2 n^2 + a_1 n + a_0$ and
$b(n) = b_1 n + b_0$, we set
\[
  r \;:=\; \frac{a_2}{b_1^2}
\]
and call this ratio the \emph{Worpitzky parameter} of $K$.

The companion note~\cite{TRINote} announces the \emph{Transcendental
Ratio Identity}: every degree-$(2,1)$ integer PCF whose limit lies in
the transcendental stratum (\textsection\ref{sec:background}) and
whose Worpitzky parameter is in the negative interior $(-1/4,0)$
satisfies $r=-2/9$.  The conjecture is supported by zero
counterexamples in approximately $585{,}000$ integer families surveyed
across the F$(2,4)$ and F$(2,5)$ complete censuses and several
targeted resonance sweeps at the next admissible rational-indicial
ratios; a single Log-stratum anomaly at $r=-1/36$ (two families,
$b_1=\pm 6$) is documented and falls outside the conjecture's hypothesis.

The present paper is the long companion of~\cite{TRINote}.  It
extends the announcement in two structurally distinct directions:

\begin{itemize}[leftmargin=1.4em]
\item \textbf{An algebraic signature
(\textsection\ref{sec:signature}).}  Under the alternative
minus-sign recurrence form
$y_{n+1} = b(n)y_n - a(n)y_{n-1}$, the
characteristic equation is $\mu^2 - b_1\mu + a_2 = 0$, with roots
$\mu_\pm = b_1(3\pm\sqrt{17})/6$ on the Trans locus $r=-2/9$.  We
prove (Theorem~\ref{thm:signature}) that, for $a_2<0$ and
$b_1\ne 0$,
\[
  \frac{a_2}{b_1^2} \,=\, -\frac{2}{9}
  \quad\Longleftrightarrow\quad
  \left|\frac{\mu_+}{\mu_-}\right| \,=\, \frac{13 + 3\sqrt{17}}{4}\,,
\]
and verify the identity by PSLQ at $\dps = 150$ with the phantom-hit
rule applied (Corollary~\ref{thm:pslq}).  This is the first explicit
algebraic identification of a quadratic irrationality (the discriminant
$17$) as an invariant of the Trans property of integer
degree-$(2,1)$ PCFs.
\item \textbf{Three negative results bounding the analytic landscape
(\textsection\ref{sec:negative}).}  We carry out, in turn, the
indicial-exponent forcing (\textsection\ref{sec:indicial}), the
leading-order Birkhoff invariant search
(\textsection\ref{sec:leading}), and the sub-leading $[n^{-1}]$
balance (\textsection\ref{sec:subleading}).  All three return
\emph{negative} answers in a precise sense: no constraint accessible
to these standard methods singles out the Trans locus among the
admissible rational-indicial loci.  The mechanism that produces the
empirical universality of $-2/9$ must therefore be \emph{global}
rather than local-asymptotic.
\end{itemize}

The two axes are not independent: $\sqrt{17}$ appears in three guises
across the paper --- as the discriminant of the characteristic
equation on the Trans locus
(\textsection\ref{sec:signature}), as the irrational ratio
$b_0/b_1=(27\pm\sqrt{17})/18$ that obstructs the integer realisation
of the indicial pair $\{1/3,2/3\}$ in minus-sign form
(\textsection\ref{sec:indicial}), and as a coefficient
$\mp\sqrt{17}\,b_1/3$ in the regular sub-leading balance on the Trans
locus (\textsection\ref{sec:subleading}).  The recurrence of
$\sqrt{17}$ in distinct asymptotic settings is, we argue, the
algebraic shadow of the global mechanism that the paper otherwise
fails to capture.

\paragraph{Organisation.}
\textsection\ref{sec:background} reviews degree-$(2,1)$ PCFs,
fixes notation, and contrasts the two recurrence sign conventions
in use.  \textsection\ref{sec:conjecture} restates the Transcendental
Ratio Identity and summarises the empirical evidence
from~\cite{TRINote} (~585\,000 families, the F$(2,5)$ census, the
resonance sweeps at $|b_1|\le 12$, and the Log anomaly at $r=-1/36$).
\textsection\ref{sec:signature} proves the
$|\mu_+/\mu_-|$-modulus identity and gives the PSLQ certificate.
\textsection\ref{sec:negative} contains the three negative results.
\textsection\ref{sec:discussion} discusses the global / local
dichotomy, the role of $\sqrt{17}$, and four open problems.
Appendix~\ref{app:methods} fixes the PSLQ methodology and the
phantom-hit rule.  Appendix~\ref{app:hashes} catalogues the SHA-256
hashes that anchor the AEAL claim manifest.

%% ==================================================================
\section{Background and notation}\label{sec:background}
%% ==================================================================

\subsection{Degree-$(2,1)$ PCFs and their strata}

Throughout, $K(a,b)$ denotes a PCF of the form~\eqref{eq:pcf} with
$a(n) = a_2 n^2 + a_1 n + a_0$, $b(n) = b_1 n + b_0$, and
$(a_0,a_1,a_2,b_0,b_1) \in \Z^5$ with $a_2 b_1 \ne 0$.  We write
$\FF(2,D) := \{(a_0,a_1,a_2,b_0,b_1)\in\Z^5 : |a_i|,|b_j|\le D,\;
a_2 b_1\ne 0\}$ for the parameter cube of side $D$.

The arithmetic strata used in this paper follow the conventions of
\cite{TRINote, F25Census}.  A convergent PCF $K$ with limit $L$ is
classified as
\(
\Rat,\,\Alg,\,\Log,\,\Trans,\,\Des
\)
according as $L \in \Q$; $L$ is algebraic of degree $\ge 2$; $L$ is a
$\Q$-linear combination of logarithms of algebraic numbers; $L$ is
transcendental and not in the previous classes (typically a M\"obius
transform of $\pi$ or an Apéry-like series); or the partial fractions
fail to converge to a finite limit.  PSLQ-based stratum identification
is performed at $\dps = 100$ with $1500$ partial quotients and a
between-truncant tolerance of $10^{-8}$, calibrated against a
$\dps = 150$ deep-validation pass.

\subsection{Two recurrence conventions}\label{sec:conventions}

The companion note~\cite{TRINote} uses the
\emph{plus-form} recurrence
\begin{equation}\label{eq:plus}
  y_{n+1} \;=\; b(n)\,y_n \,+\, a(n)\,y_{n-1}\qquad(\text{plus-form})
\end{equation}
which is standard for the convergent sequences $P_n,Q_n$ of a CF
written as in~\eqref{eq:pcf}.  The asymptotic analysis of
\textsection\ref{sec:negative} below uses instead the
\emph{minus-form} recurrence
\begin{equation}\label{eq:minus}
  y_{n+1} \;=\; b(n)\,y_n \,-\, a(n)\,y_{n-1}\qquad(\text{minus-form})
\end{equation}
which arises when the sign of $a$ is absorbed into the recurrence
relation rather than into the coefficient.  The two are conjugate by
$y_n \mapsto (-1)^n y_n$ \emph{up to} the sign of $a(n)$; the locus
$r = a_2/b_1^2$ is invariant, but the indicial exponents and
characteristic roots transform.  Specifically, on the locus
$r = -2/9$:
\begin{itemize}[leftmargin=1.4em]
\item The plus-form indicial equation
$\lambda^2 - \lambda - r = 0$ has rational roots $\{1/3,2/3\}$
(Lemma~3.2 of~\cite{TRINote}, restated here as
Theorem~\ref{thm:plus-indicial}).
\item The minus-form characteristic equation $\mu^2 - b_1\mu + a_2=0$
has irrational roots $\mu_\pm = b_1(3\pm\sqrt{17})/6$
(Lemma~\ref{lem:minus-roots}).
\end{itemize}
We use the plus-form for the definitional content of
\textsection\ref{sec:conjecture} and the minus-form for the
asymptotic analysis of \textsection\textsection\ref{sec:signature}--\ref{sec:negative}.
A summary correspondence appears in Table~\ref{tab:conventions}.

\begin{table}[t]
\centering
\caption{Sign-convention dictionary on the Trans locus $r=-2/9$.}
\label{tab:conventions}
\begin{tabular}{lll}
\toprule
Form & Recurrence & Eigenstructure on $r=-2/9$ \\
\midrule
plus & $y_{n+1} = b(n)y_n + a(n)y_{n-1}$
& indicial pair $\{1/3,2/3\}\subset\Q$ \\
minus & $y_{n+1} = b(n)y_n - a(n)y_{n-1}$
& char.\ roots $\mu_\pm = b_1(3\pm\sqrt{17})/6\notin\Q$ \\
\bottomrule
\end{tabular}
\end{table}

\subsection{PSLQ and the phantom-hit rule}\label{sec:pslq}

PSLQ integer-relation detection is performed using
\texttt{mpmath.pslq} (Python~3, mpmath~1.3.0).  Working precision is
$\dps = 100$ for census passes, $\dps = 150$ for the algebraic
signature certificate of \textsection\ref{sec:signature}, and
$\dps = 200$--$300$ for the sub-leading verification of
\textsection\ref{sec:subleading}.  Tolerance is fixed at $10^{-50}$ and
the integer-coefficient bound is set per task as documented in
Appendix~\ref{app:methods}.

\paragraph{Phantom-hit rule.}
Following the convention of the T2A relay
pipeline~\cite{TRINote}, every PSLQ output relation is screened for a
\emph{phantom hit}: if the integer coefficient on the target value
(the quantity whose closed form is being tested) is zero, the relation
is discarded as a phantom of the search basis.  In the T2A pipeline,
$38$ candidate relations were rejected by this rule.  The certificates
recorded in this paper are explicitly labelled with the target index
and the relevant non-zero coefficient.

%% ==================================================================
\section{The conjecture and empirical evidence}\label{sec:conjecture}
%% ==================================================================

We restate the central conjecture and summarise the empirical evidence
from~\cite{TRINote}; the reader familiar with the announcement note
may skip to \textsection\ref{sec:signature}.

\subsection{The plus-form indicial equation and the conjecture}

\begin{theorem}[Plus-form indicial pair on $r=-2/9$]
\label{thm:plus-indicial}
Let $K(a,b)$ be a degree-$(2,1)$ PCF satisfying $a_2/b_1^2 = -2/9$.
Then the plus-form recurrence~\eqref{eq:plus} has indicial equation
$\lambda^2 - \lambda - r = 0$ with rational root pair
$\{\lambda_-,\lambda_+\} = \{1/3,2/3\}$.
\end{theorem}
\begin{proof}
With $r = -2/9$, the indicial equation is $\lambda^2 - \lambda + 2/9 =
0$, whose discriminant is $1 - 8/9 = 1/9$, giving roots
$\lambda = (1 \pm 1/3)/2 \in \{1/3, 2/3\}$.  See \cite[Lemma~2.2]%
{TRINote} for the derivation of the indicial equation itself.
\end{proof}

\begin{conjecture}[Transcendental Ratio Identity, \cite{TRINote}]
\label{conj:tri}
Let $K(a,b)$ be a degree-$(2,1)$ PCF with integer coefficients,
$a_2,b_1 \ne 0$, converging to a limit $L$ in the Trans stratum and
with Worpitzky parameter $r = a_2/b_1^2$ in the negative interior
$(-1/4,0)$.  Then
\[
  r \;=\; -\tfrac{2}{9}\,.
\]
Equivalently, the plus-form indicial exponents at infinity are
$\{1/3,2/3\}$.
\end{conjecture}

\subsection{Empirical evidence}\label{sec:evidence}

Table~\ref{tab:censuses} summarises the four computational
experiments of~\cite{TRINote}; Table~\ref{tab:resonance} summarises
the resonance sweeps; Table~\ref{tab:log-anomaly} catalogues the Log
anomaly.  No interior-stratum Trans counterexample to
Conjecture~\ref{conj:tri} appeared in any of the
$\sim\!585{,}000$ integer families surveyed.

\begin{table}[t]
\centering
\caption{Stratum counts in $\FF(2,4)$ and $\FF(2,5)$
(reproduced from \cite{TRINote}).}
\label{tab:censuses}
\begin{tabular}{lrrl}
\toprule
Search & Trans & Log & Worpitzky parameter $r$ \\
\midrule
$\FF(2,4)$ ($D=4$, complete)
  & $24$ & $0$  & $-2/9$ \\
$\FF(2,5)$ ($D=5$, neg.\ interior)
  & $56$ & $12$ & $-2/9$ \\
$\FF(2,5)$ ($D=5$, Brouncker boundary)
  & $14$ & $0$  & $+1/4$ \\
\bottomrule
\end{tabular}
\end{table}

\begin{table}[t]
\centering
\caption{Targeted resonance searches at admissible rational-indicial
ratios in the negative Worpitzky interior. All runs at $\dps=100$.
Reproduced from~\cite{TRINote}.}
\label{tab:resonance}
\begin{tabular}{lrrr}
\toprule
Scope & Convergent families & Trans & Log \\
\midrule
$b_1\in\{\pm4,\pm5\}$, $r=-3/16$ & $2{,}384$ & $0$ & $0$ \\
$b_1\in\{\pm4,\pm5\}$, $r=-4/25$ & $2{,}396$ & $0$ & $0$ \\
$b_1\in\{\pm4,\pm5\}$, $r=-6/25$ & $2{,}394$ & $0$ & $0$ \\
\midrule
$b_1\in\{\pm6,\pm7\}$ (full sweep) & $175{,}686$ & $0$ & $2$ \\
$b_1\in\{\pm8,\dots,\pm12\}$ (motivated targets)
                                 & $259{,}280$ & $0$ & $0$ \\
\midrule
Combined                          & $442{,}140$ & $0$ & $2$ \\
\bottomrule
\end{tabular}
\end{table}

\begin{table}[t]
\centering
\caption{Log-stratum anomaly at $r=-1/36$
(both families have limit $L=\pm 2/\log 2$, $\dps=150$ residual
$5.65\!\times\!10^{-150}$).}
\label{tab:log-anomaly}
\begin{tabular}{ccccc}
\toprule
$a_2$ & $a_1$ & $a_0$ & $b_1$ & $b_0$ \\
\midrule
$-1$ & $0$ & $0$ & $\phantom{-}6$ & $\phantom{-}3$ \\
$-1$ & $0$ & $0$ & $-6$ & $-3$ \\
\bottomrule
\end{tabular}
\end{table}

\begin{remark}[Phantom-hit discipline]\label{rem:phantom-discipline}
The T2A pipeline used to generate Tables~\ref{tab:censuses}--\ref{tab:resonance}
applies the phantom-hit rule of \textsection\ref{sec:pslq} at every
PSLQ stage: $38$ candidate relations were rejected during the
classification.  The two known phantom traps for the bases used here
are $\zeta(2) = \pi^2/6$ and the golden ratio $\varphi=(1+\sqrt{5})/2$,
arising whenever $\pi$ or $\sqrt 5$ enters the basis; neither affected
any Trans or Log classification reported above.
\end{remark}

\subsection{The Brouncker boundary class}

The $14$ families at $r = +1/4$ in $\FF(2,5)$ lie on the positive
Worpitzky boundary $\Delta = 2$ (where $\Delta = 1+4r$ is the
plus-form indicial discriminant) and converge logarithmically
to M\"obius transforms of $4/\pi$.  They are excluded from
Conjecture~\ref{conj:tri} by the negative-interior hypothesis
$r\in(-1/4,0)$, and do not contradict the Trans Identity.  See
\cite[Theorem~4.1]{TRINote}.

%% ==================================================================
\section{An algebraic signature: the
\texorpdfstring{$|\mu_+/\mu_-|$}{|mu+/mu-|} identity}
\label{sec:signature}
%% ==================================================================

We turn to the new content.  Throughout this section we work in the
\emph{minus-form} recurrence~\eqref{eq:minus}.

\subsection{The minus-form characteristic equation}

\begin{lemma}[Minus-form characteristic roots on $r=-2/9$]
\label{lem:minus-roots}
Let $K(a,b)$ be a degree-$(2,1)$ PCF, regarded as a minus-form
recurrence $y_{n+1} = b(n)y_n - a(n)y_{n-1}$.  The Birkhoff ansatz at
infinity $y_n \sim \Gamma(n+1)\,\mu^n n^\alpha$
(\textsection\ref{sec:indicial} below) gives the characteristic
equation
\begin{equation}\label{eq:minus-char}
  \mu^2 \,-\, b_1\,\mu \,+\, a_2 \;=\; 0
\end{equation}
with roots $\mu_\pm = \tfrac{1}{2}\bigl(b_1 \pm \sqrt{b_1^2 - 4
a_2}\bigr)$.  In particular, on the locus $a_2/b_1^2 = -2/9$ with
$a_2 < 0$ and $b_1 \ne 0$,
\[
  \mu_\pm \;=\; \frac{b_1\,(3 \pm \sqrt{17})}{6}\,,
\]
both irrational quadratic over $\Q$.
\end{lemma}
\begin{proof}
Substituting $y_n = \Gamma(n+1)\mu^n n^\alpha$ into~\eqref{eq:minus},
expanding $\Gamma(n+2)/\Gamma(n+1) = n+1$ and $(1\pm 1/n)^\alpha
= 1\pm\alpha/n + O(1/n^2)$, multiplying through by $n^{-2}$, and
reading the leading $[n^1]$ coefficient (after cancellation against
the $\Gamma$-recursion) gives the relation $b_1\mu = \mu^2 + a_2$,
i.e.~\eqref{eq:minus-char}.  See \cite[\textsection 2]{SprintW1Report}
for the full sympy verification.  Substituting $a_2 = -2b_1^2/9$ into
the discriminant gives $b_1^2 - 4 a_2 = b_1^2(1 + 8/9) = 17 b_1^2/9$
and $\sqrt{b_1^2-4a_2} = |b_1|\sqrt{17}/3$.
\end{proof}

\begin{theorem}[Algebraic signature of the Trans locus]
\label{thm:signature}
Let $K(a,b)$ be a degree-$(2,1)$ PCF with $a_2 < 0$ and $b_1 \ne 0$,
and let $\mu_\pm$ be the minus-form characteristic roots
of~\eqref{eq:minus-char}.  Set $\rho := |\mu_+/\mu_-|$.  Then
\begin{equation}\label{eq:signature}
  \frac{a_2}{b_1^2} \;=\; -\frac{2}{9}
  \quad\Longleftrightarrow\quad
  \rho \;=\; \frac{13 + 3\sqrt{17}}{4}\,.
\end{equation}
\end{theorem}
\begin{proof}
Set $q := \mu_+/\mu_-$ (so $\rho = |q|$).  Vieta's formulas applied
to~\eqref{eq:minus-char} give $\mu_+ + \mu_- = b_1$ and
$\mu_+\mu_- = a_2$, whence
\[
\frac{a_2}{b_1^2}
\;=\; \frac{\mu_+\mu_-}{(\mu_+ + \mu_-)^2}
\;=\; \frac{q}{(1+q)^2}\,.
\]
($\Rightarrow$) Substituting $a_2/b_1^2 = -2/9$ gives the quadratic
$2q^2 + 13 q + 2 = 0$, with roots
\(q = (-13 \pm 3\sqrt{17})/4\).
Both roots are negative; their moduli are
\(|q| = (13 + 3\sqrt{17})/4\) (for the root with $|q|>1$) and
\(|q| = (13 - 3\sqrt{17})/4 = 4/(13+3\sqrt{17})\) (the reciprocal).
Either choice of labelling $\mu_+$ versus $\mu_-$ yields
$\rho = (13+3\sqrt{17})/4$ after taking the dominant root in absolute
value.
\\
($\Leftarrow$) Conversely, if $\rho = (13+3\sqrt{17})/4$, then
$q \in \{-(13+3\sqrt{17})/4,\ +(13+3\sqrt{17})/4\}$ (two
sign choices for the quotient).  Substituting into
$a_2/b_1^2 = q/(1+q)^2$:
\begin{itemize}[leftmargin=1.4em]
\item $q = -(13+3\sqrt{17})/4$ gives $a_2/b_1^2 = -2/9$;
\item $q = +(13+3\sqrt{17})/4$ gives $a_2/b_1^2 = +2/17$
(positive Worpitzky region).
\end{itemize}
The hypothesis $a_2 < 0$ excludes the second branch.
\end{proof}

\begin{remark}[Vieta-tautological character and its content]
\label{rem:vieta-tautology}
Theorem~\ref{thm:signature} is, formally, a Vieta tautology: the
ratio $|\mu_+/\mu_-|$ is a function of $a_2/b_1^2$ alone.  What is
\emph{non-tautological} is that the specific irrational
$(13+3\sqrt{17})/4$ on the right-hand side displays the discriminant
$17$ explicitly.  As we shall see in
\textsection\textsection\ref{sec:indicial}--\ref{sec:subleading},
$\sqrt{17}$ recurs in two further asymptotic settings, and is the
algebraic shadow of the Trans property under all three local probes
that fail to characterise it.
\end{remark}

\subsection{PSLQ certificate}

\begin{corollary}[PSLQ certificate]\label{thm:pslq}
At working precision $\dps = 150$ and tolerance $10^{-50}$:
\begin{enumerate}[label={\upshape(\arabic*)}]
\item PSLQ on the basis
$\{1,\sqrt{17},\rho\}$ with $\rho = |\mu_+/\mu_-|$ on a
representative Trans family returns the integer relation
$[13,3,-4]$, i.e.\ $13 + 3\sqrt{17} - 4\rho = 0$.  The target
coefficient (on $\rho$) is $-4 \ne 0$; the phantom-hit rule
passes.  Reconstruction error vanishes to $> 100$ decimal digits.
\item PSLQ on the wider basis
$\{1,\sqrt{2},\sqrt{3},\sqrt{5},\sqrt{17},\pi,\log 2,\rho\}$
returns $[-13,0,0,0,-3,0,0,4]$, the same identity with
$\rho$-coefficient $4 \ne 0$.
\end{enumerate}
This follows immediately from Theorem~\ref{thm:signature}, with the
relation confirmed numerically by PSLQ at \texttt{dps}=150.
\end{corollary}
\begin{proof}[Computational verification]
Script \texttt{forcing\_condition.py} of session
\texttt{2026-04-27/SPRINT-W2-TRANS-INDICIAL-SURVEY/sprint\_w2/}.  See
Appendix~\ref{app:hashes} for the SHA-256 hash of the run log.
\end{proof}

%% ==================================================================
\section{Three negative results}\label{sec:negative}
%% ==================================================================

We now record three negative results.  In each, a standard local
asymptotic method --- (i) indicial-exponent forcing, (ii) leading-order
Birkhoff invariant search, (iii) sub-leading $[n^{-1}]$ resonance ---
is shown to fail to single out the Trans locus among admissible
rational-indicial loci.

\subsection{Indicial exponent forcing under the minus form}
\label{sec:indicial}

The W1 sprint of this paper~\cite{SprintW1Report} works in the
minus-form~\eqref{eq:minus} and asks: which integer parameter tuples
$(a_0,a_1,a_2,b_0,b_1)$ realise the indicial pair $\{1/3,2/3\}$ at
infinity \emph{in minus-form}?  Note this is a different question
from Theorem~\ref{thm:plus-indicial}: in plus-form, $\{1/3,2/3\}$ is
a \emph{consequence} of $r=-2/9$; in minus-form, the corresponding
pair on the Trans locus is irrational by Lemma~\ref{lem:minus-roots},
and asking for the rational pair $\{1/3,2/3\}$ is a non-trivial
constraint.

\begin{proposition}[Minus-form indicial closed form,
\cite{SprintW1Report}]
\label{prop:minus-indicial}
Under the Birkhoff ansatz $y_n \sim \Gamma(n+1)\,\mu^n n^\alpha$
applied to~\eqref{eq:minus}, the next-to-leading $[n^0]$ balance,
after using $\mu^2 = b_1\mu - a_2$, has the unique solution
\begin{equation}\label{eq:alpha-closedform}
  \alpha(\mu) \;=\; -\,\frac{(b_1 - b_0)\mu \;+\; (a_1 - a_2)}%
                            {b_1\mu - 2 a_2}\,.
\end{equation}
Each minus-form characteristic root $\mu_\pm$ produces an indicial
exponent $\alpha_\pm := \alpha(\mu_\pm)$.
\end{proposition}
\begin{proof}
Direct sympy substitution and verification; see
\cite[\textsection 3]{SprintW1Report}, script
\texttt{indicial\_symbolic.py}.
\end{proof}

\begin{lemma}[Compatibility ideal for the pair $\{1/3,2/3\}$]
\label{lem:compat-ideal}
The pair $\{\alpha_-,\alpha_+\} = \{1/3,2/3\}$, i.e.\ the
sum-and-product conditions $S = \alpha_++\alpha_- = 1$ and $P =
\alpha_+\alpha_- = 2/9$, is equivalent to the system
\[
  \begin{cases}
    a_1 - 2 a_2 = 0,\\
    a_2 \;+\; 9 b_0^2 - 27 b_0 b_1 + 20 b_1^2 \;=\; 0,
  \end{cases}
\]
defining the compatibility ideal $C \subset
\Q[a_0,a_1,a_2,b_0,b_1]$.
\end{lemma}
\begin{proof}
Substitute~\eqref{eq:alpha-closedform} into the symmetric functions
$S$ and $P$ and simplify using Vieta on $(\mu_+,\mu_-)$;
see~\cite[\textsection 4]{SprintW1Report}.
\end{proof}

\begin{theorem}[Integer-grid obstruction, negative result]
\label{thm:integer-obstruction}
Fix $a_2/b_1^2 = -2/9$ and set $\varrho := b_0/b_1$.  The
intersection of the Trans locus with the compatibility ideal
$C$ of Lemma~\ref{lem:compat-ideal} projects onto the constraint
\[
  -9\,\varrho^2 + 27\,\varrho - 20 \;=\; \frac{a_2}{b_1^2}
  \;=\; -\frac{2}{9}\,,
\]
i.e.\ $81\,\varrho^2 - 243\,\varrho + 178 = 0$, with roots
$\varrho = (27 \pm \sqrt{17})/18$, both irrational.  Consequently no
integer-coefficient family realises both the Trans locus
$a_2/b_1^2 = -2/9$ and the minus-form indicial pair $\{1/3,2/3\}$
simultaneously.  Direct enumeration on the integer grid
$|b_1|\le 30$, $|b_0|\le 30$, $|a_2|\le 200$
($\sim 7.3\times 10^5$ triples) confirms zero simultaneous
solutions.
\end{theorem}
\begin{proof}
The root computation is elementary; the enumeration is the script
\texttt{compatibility.py} of \cite{SprintW1Report}.
\end{proof}

\begin{remark}[Reconciliation with Theorem~\ref{thm:plus-indicial}]
Theorem~\ref{thm:integer-obstruction} does not contradict
Theorem~\ref{thm:plus-indicial}: the two work in opposite sign
conventions and produce consistent but distinct rational pictures.
The conjecture~\ref{conj:tri} as stated uses the plus-form, where
indicial $\{1/3,2/3\}$ on the Trans locus is automatic.
Theorem~\ref{thm:integer-obstruction} says, instead, that under the
minus-form (which is the natural setting for Birkhoff
asymptotics with the $\Gamma(n+1)$ leading factor), the candidate
rational indicial pair $\{1/3,2/3\}$ requires irrational coefficient
ratios involving $\sqrt{17}$, and hence is unrealisable over $\Z$.
The integer-Trans empirical content of $-2/9$ \emph{is} reflected in
the minus-form picture, but only through the algebraic signature
$\sqrt{17}$, not as a rational indicial pair.
\end{remark}

\subsection{Leading-order Birkhoff invariant search}\label{sec:leading}

The W2 sprint~\cite{SprintW2Report} surveys $50$ integer Trans
families and $20$ non-Trans control families, computing for each
the minus-form characteristic data
$(\mu_\pm, |\mu_+/\mu_-|, \alpha_\pm, S, P, D)$, where
$D = (\alpha_+-\alpha_-)^2$ is the indicial discriminant.

\begin{observation}[No new leading-order invariant, negative result]
\label{prop:leading-invariant}
Among the leading two-term Birkhoff invariants
$(|\mu_+/\mu_-|, S, P, D)$ on the integer Trans locus
$a_2/b_1^2 = -2/9$:
\begin{itemize}[leftmargin=1.4em]
\item $|\mu_+/\mu_-|$ is constant $= (13+3\sqrt{17})/4$
(Theorem~\ref{thm:signature}, Vieta-tautological);
\item $S$ takes three distinct rational values across the survey
($\{-3,-1,+1\}$);
\item $P$ takes nine distinct values in $[-2.235,+2.235]$;
\item $D$ takes eight distinct values in $[0.059,9.94]$.
\end{itemize}
Consequently, no rational-coefficient invariant accessible to the
leading two-term Birkhoff balance, beyond the Vieta-tautological
$|\mu_+/\mu_-|$ identity, distinguishes the Trans locus.
\end{observation}
\begin{proof}
Computational survey, script \texttt{indicial\_survey.py} of
\cite{SprintW2Report}; the table of values and the variability ranges
are reproduced verbatim from the sprint report.
\end{proof}

\begin{remark}[Class separation by $r$]
The ratio $|\mu_+/\mu_-|$ is, by Vieta, a function of $r=a_2/b_1^2$
alone; it therefore separates the strata Trans, Log, Alg by
construction (Trans has the unique value $(13+3\sqrt{17})/4$; Log
families take five distinct values; Alg degenerate or complex).
This separation is built into the definition of the strata.  The
content of Observation~\ref{prop:leading-invariant} is that
\emph{nothing else} at leading two-term order is constant across
Trans.
\end{remark}

\subsection{The sub-leading $[n^{-1}]$ balance}\label{sec:subleading}

The W3 sprint~\cite{SprintW3Report} extends the Birkhoff ansatz to
\[
  y_n \;\sim\; \Gamma(n+1)\,\mu^n\,n^\alpha\,
  \left(1 + \frac{c_1}{n} + \frac{c_2}{n^2} + \cdots\right)
\]
and derives the $[n^{-1}]$ balance, linear in $c_1$ and $c_2$.  The
$c_2$-coefficient of the balance equals $\mu^2 - b_1\mu + a_2$, which
vanishes on the characteristic locus by Lemma~\ref{lem:minus-roots};
$c_2$ is therefore not constrained at this order.  The
$c_1$-coefficient $A$, after using $\mu^2 = b_1\mu - a_2$ and
substituting the closed form~\eqref{eq:alpha-closedform} for $\alpha$,
factors as follows.

\begin{proposition}[Sub-leading $c_1$ coefficient, \cite{SprintW3Report}]
\label{prop:c1-coefficient}
With notation as above, the coefficient of $c_1$ at the $[n^{-1}]$
order of the minus-form Birkhoff balance is, after reduction,
\begin{equation}\label{eq:c1-A}
  A \;=\; \frac{(a_2 - b_1\mu)(4 a_2 - b_1^2)}%
                {a_2 b_1 + 2 a_2\mu - b_1^2\mu}\,.
\end{equation}
\end{proposition}
\begin{proof}
Symbolic computation in sympy; see~\cite[\textsection 2]{SprintW3Report}
and script \texttt{subleading\_birkhoff.py}.  The numerator
$A_{\text{num}} = (a_2-b_1\mu)(4 a_2 - b_1^2)$ is verified directly;
the denominator $A_{\text{den}}$ is an explicit rational
expression in $(a_2,b_1,\mu)$ and is non-zero on the Trans locus.
\end{proof}

\begin{corollary}[No sub-leading resonance on the Trans locus,
negative result]
\label{cor:no-resonance}
$A_{\text{num}}$ vanishes if and only if (i) $a_2 = b_1\mu$, which
combined with the characteristic equation~\eqref{eq:minus-char}
forces $\mu = 0$ and is therefore degenerate; or (ii)
$4 a_2 = b_1^2$, which is the algebraic-stratum (Alg)
discriminant-zero locus.  Neither contains the Trans locus
$a_2/b_1^2 = -2/9$.  On the Trans locus,
\[
  A\bigl|_{\mu = \mu_\pm}
  \;=\; \mp\,\frac{\sqrt{17}\,b_1}{3}
  \;\neq\;0\,,
\]
so the $[n^{-1}]$ balance uniquely determines $c_1$ as $-B/A$
(with $B$ the $c_1$-and-$c_2$-free remainder), and yields no new
constraint on the parameters.
\end{corollary}
\begin{proof}
Combine Proposition~\ref{prop:c1-coefficient} with
Lemma~\ref{lem:minus-roots} for the explicit values of $A$ at
$\mu_\pm$.
\end{proof}

\begin{observation}[Convergence-rate experiment, negative result]
\label{prop:rate}
For the minus-form recurrence on a representative panel of $10$
Trans, $5$ Log, and $5$ Alg integer families at $\mathrm{mp.dps}=200$
and $N=600$ iterations, the partial-fraction error
$|r_n - L|$ (with $r_n = p_n/q_n$ and $L$ the high-$n$ limit) reaches
$10^{-200}$ at $n\approx 300$ for every Trans and Log family;
the classical normalisation $\rho_n := |r_n-L|\cdot q_n^2$ explodes
($\sim\!10^{2620}$ at $n=500$) because $q_n \sim \Gamma(n+1)\mu_+^n$;
the natural rate constant
\[
  \lim_{n\to\infty}\frac{|r_n - L|\cdot q_n\,q_{n+1}}{|\mu_-/\mu_+|^n}
\]
exists but reduces, on the Trans locus, to the
Vieta-tautological quantity
$|\mu_-/\mu_+| = 4/(13+3\sqrt{17}) = (13-3\sqrt{17})/4$.  No new
amplitude or coefficient invariant is exposed by the convergence-rate
experiment at this precision.
\end{observation}
\begin{proof}
Script \texttt{convergence\_rate.py} of~\cite{SprintW3Report}.
The hyperexponential rate is a direct consequence of the
$\Gamma(n+1)$ factor in the dominant Birkhoff solution; the rate
\emph{constant} after $\Gamma$-removal is the characteristic-root ratio.
\end{proof}

\begin{observation}[Direct PSLQ on the Trans limits, negative result]
\label{prop:pslq-direct}
At $\mathrm{mp.dps} = 300$, $N=500$, and integer-coefficient bound
$10^4$:
\begin{enumerate}[label={\upshape(\roman*)}]
\item PSLQ on $\{1,L,L^2,L^3,L^4\}$ for each of $10$ representative
Trans CF limits $L$ returns no integer relation.
\item Pairwise PSLQ on $\{1,L_i,L_j,L_iL_j,L_i^2,L_j^2\}$ returns ten
relations among the $45$ pairs; all ten are tautological consequences
of the linear scaling $L(\lambda b_1) = \lambda L(b_1)$ on the
sub-family with $(a_0,a_1,b_0) = (0,0,0)$, and are not novel
algebraic relations between distinct CFs.
\item PSLQ on $\{1,L,\sqrt{17},\pi,\log 2\}$ for each of the $10$
Trans limits returns no integer relation.
\end{enumerate}
\end{observation}
\begin{proof}
Script \texttt{direct\_algebraic.py} of~\cite{SprintW3Report}.
The phantom-hit rule is applied at every step; the scale-invariance
hits in (ii) are not phantom in the strict sense (the
$L$-coefficient is non-zero) but are recognised tautological by post
hoc inspection.
\end{proof}

%% ==================================================================
\section{Discussion and open problems}\label{sec:discussion}
%% ==================================================================

\subsection{Local versus global mechanism}

The combined effect of \textsection\textsection\ref{sec:indicial}--\ref{sec:subleading}
is that no \emph{local} asymptotic method --- characteristic
exponents, leading invariants, sub-leading resonance, convergence
rate --- distinguishes the Trans locus from other admissible
rational-indicial loci to within the precision and search bounds
employed.  The empirical universality of $r=-2/9$ over
$\sim\!585{,}000$ integer families is therefore very probably the
shadow of a \emph{global} structural property, of a kind not visible
in expansions about $n=\infty$ alone.

We see two natural candidate frameworks for capturing this global
property.

\subsubsection{Galois theory of linear difference equations}
The minus-form recurrence~\eqref{eq:minus} is a linear difference
operator $\Delta = E^2 - b(n) E + a(n)$ on the field $\Q(n)$.  The
Galois group of the splitting extension of $\Delta$ over $\Q(n)$
(in the sense of Picard-Vessiot for difference equations) classifies
the formal solutions modulo rational equivalences.  Whether the
\emph{integer-realisability} condition (i.e.\ that some pair of
solutions of $\Delta$ has a transcendental ratio of integer
convergents) corresponds to a specific Galois group --- and whether
that group is detected by $r=-2/9$ --- is to our knowledge open.
The recurrence of $\sqrt{17}$ across our three local probes
suggests that any such Galois-theoretic answer must involve the
quadratic field $\Q(\sqrt{17})$.

\subsubsection{Isomonodromic deformation theory}
The plus-form indicial pair $\{1/3,2/3\}$ matches the Schwarz-list
dihedral entry with exponent difference $1/3$~\cite{Schwarz}.  The
companion note~\cite{TRINote} speculates on a Schwarz-monodromic
mechanism behind the conjecture; in the present paper we strengthen
this speculation by noting that the irrationality $\sqrt{17}$ which
shadows the Trans property in three asymptotic settings is exactly
the discriminant of the minus-form characteristic polynomial on the
Trans locus, suggesting a quadratic-field constraint (rather than
the rational $\{1/3,2/3\}$) is the ``correct'' invariant for
isomonodromic considerations.  A Painlevé-VI-type interpretation,
in which the Trans locus appears as a special isomonodromic stratum
of a difference-equation analogue of Painlevé~VI with parameters on
the Schwarz dihedral, is a natural structural problem
(Problem~\ref{prob:painleve} below).

\subsection{The Log anomaly at $r=-1/36$}

The two Log families at $b_1=\pm 6$, $r=-1/36$ (Table~\ref{tab:log-anomaly})
sit at the intersection of (i) the Bauer--Stern doubling parameter
of the Log stratum and (ii) the rational-indicial admissibility set
$\{1/q,(q-1)/q\}$ for $q=6$.  Whether this is the unique boundary
phenomenon, or the tip of a finer Log-stratum classification, is
open (Problem~\ref{prob:log-anomaly}).

\subsection{Open problems}

\begin{problem}[Galois-of-difference-equations form]
\label{prob:galois}
Define the Picard-Vessiot Galois group $G_K$ of a degree-$(2,1)$
integer PCF.  Is there a property of $G_K$ that characterises the
Trans locus?  Equivalently: is the empirical universality of
$r=-2/9$ a constraint on the conjugacy class of $G_K$ (e.g.\ a
specific finite quotient)?  We expect any such characterisation to
factor through the quadratic field $\Q(\sqrt{17})$.
\end{problem}

\begin{problem}[Isomonodromic / Painlevé form]\label{prob:painleve}
Construct (or rule out) a difference-equation analogue of
Painlevé~VI in which the Trans locus appears as a special
isomonodromic stratum at parameters on the Schwarz-dihedral
$\{1/3,2/3\}$.  A natural starting point is the Jimbo-Sakai
$q$-analogue.
\end{problem}

\begin{problem}[Rational-coefficient extension]
\label{prob:rational}
Does the Transcendental Ratio Identity (Conjecture~\ref{conj:tri})
survive when the integer hypothesis on $(a_0,a_1,a_2,b_0,b_1)$ is
relaxed to rational?
\end{problem}

\begin{problem}[Log anomaly classification]\label{prob:log-anomaly}
Classify all degree-$(2,1)$ integer PCFs with $a_2/b_1^2 = -1/36$
producing Log limits.  Are the two known $b_1=\pm 6$ families
exceptional, or do larger $b_1$ produce a Log sub-family at this
boundary parameter?
\end{problem}

%% ==================================================================
\section{Conclusion}\label{sec:conclusion}
%% ==================================================================

The Transcendental Ratio Identity for degree-$(2,1)$ integer PCFs is
empirically robust ($\sim\!585{,}000$ families surveyed, zero
counterexamples), algebraically fingerprinted by the quadratic
irrationality $\sqrt{17}$, and not accessible to any of the three
standard local asymptotic methods we have examined.  The mechanism
producing the empirical universality of $r=-2/9$ must be global ---
plausibly Galois-theoretic or isomonodromic --- and lies beyond the
scope of the techniques deployed in this paper.

The role of $\sqrt{17}$ is, we believe, the most striking output of
the present investigation.  It appears as (i) the characteristic-root
discriminant of the minus-form recurrence on the Trans locus
(Lemma~\ref{lem:minus-roots}); (ii) the irrationality of the
$b_0/b_1$ ratios solving the minus-form indicial-pair compatibility
ideal (Theorem~\ref{thm:integer-obstruction}); and (iii) the value
of the regular sub-leading $c_1$-coefficient at $\mu = \mu_\pm$
(Corollary~\ref{cor:no-resonance}).  We conjecture that any future
proof of Conjecture~\ref{conj:tri} will route through the quadratic
field $\Q(\sqrt{17})$.

We invite specialists in the Galois theory of difference equations
and in isomonodromic deformation theory to consider the structural
problems posed in \textsection\ref{sec:discussion}.

%% ==================================================================
\appendix
\section{Computational methodology}\label{app:methods}
%% ==================================================================

All computations use Python~3 with \texttt{mpmath~1.3.0} and
\texttt{sympy~1.14.0}.  The PSLQ implementation is
\texttt{mpmath.pslq} with tolerance $10^{-50}$.  Working precision is
set per task:
\begin{itemize}[leftmargin=1.4em]
\item $\dps = 100$ for stratum classification across $\FF(2,4)$,
$\FF(2,5)$ and the resonance sweeps;
\item $\dps = 150$ for the algebraic signature certificate
(Corollary~\ref{thm:pslq});
\item $\dps = 200$ for the convergence-rate experiment
(Observation~\ref{prop:rate});
\item $\dps = 300$ for the direct algebraic identification of the
Trans CF limits (Observation~\ref{prop:pslq-direct}).
\end{itemize}
The integer-coefficient bound \texttt{maxcoeff} is documented per
PSLQ call; for Corollary~\ref{thm:pslq} no explicit bound was needed
(the relation has tiny coefficients), and for
Observation~\ref{prop:pslq-direct} the bound is fixed at $10^{4}$.

The phantom-hit rule of \textsection\ref{sec:pslq} is applied at
every PSLQ step.  All scripts are archived in the bridge repository
\href{https://github.com/papanokechi/siarc-relay-bridge}%
{papanokechi/siarc-relay-bridge}, and every quoted numerical claim
has an entry in the AEAL claim manifest (\texttt{claims.jsonl}) of
the relevant session, recording \texttt{evidence\_type}, \texttt{dps},
script filename, and SHA-256 hash of the run log.

%% ==================================================================
\section{SHA-256 hashes for the AEAL manifest}\label{app:hashes}
%% ==================================================================

For reproducibility, the SHA-256 hashes of the principal scripts
(at the commit underlying this paper) are tabulated in
Table~\ref{tab:hashes}.  Hashes for the W1, W2, and W3 sprints are
recorded in the corresponding session \texttt{claims.jsonl} files.

\begin{table}[t]
\centering
\caption{SHA-256 hashes of the principal scripts for the
algebraic-signature and negative-result sections.}
\label{tab:hashes}
\begin{tabular}{ll}
\toprule
Script & SHA-256 (truncated) \\
\midrule
\texttt{indicial\_symbolic.py} (W1) & see W1 \texttt{claims.jsonl} \\
\texttt{compatibility.py} (W1) & see W1 \texttt{claims.jsonl} \\
\texttt{indicial\_survey.py} (W2) & see W2 \texttt{claims.jsonl} \\
\texttt{forcing\_condition.py} (W2) & see W2 \texttt{claims.jsonl} \\
\texttt{subleading\_birkhoff.py} (W3) & \texttt{08eddada\dots dcc0b654} \\
\texttt{convergence\_rate.py} (W3) & \texttt{b7742ebb\dots d1e4ecb} \\
\texttt{direct\_algebraic.py} (W3) & \texttt{88b14129\dots 23685188f} \\
\bottomrule
\end{tabular}
\end{table}

%% ==================================================================
\section*{Computational and AI-Assistance Disclosure}
%% ==================================================================

\textbf{Reproducibility.}  All computations underlying this paper
are reproducible from publicly archived scripts and datasets.  The
classification pipelines and the PSLQ certificates are archived at
the bridge repository
\href{https://github.com/papanokechi/siarc-relay-bridge}%
{papanokechi/siarc-relay-bridge}, in the sessions referenced
in~\cite{TRINote, SprintW1Report, SprintW2Report, SprintW3Report}.
Every numerical claim is logged as a \texttt{claims.jsonl} entry
with \texttt{evidence\_type = computation}, the working precision
\texttt{dps}, the script filename, and a SHA-256 hash of the result
file.

\textbf{Phantom-hit discipline.}  All PSLQ relations reported in
this paper are accompanied by an explicit non-zero target-coefficient
check (the phantom-hit rule of \textsection\ref{sec:pslq}).  In the
T2A pipeline used for \textsection\ref{sec:conjecture}, $38$
candidate relations were rejected by this rule; the certificates
of \textsection\ref{sec:signature} pass the rule with target
coefficients $-4$ (basis 1) and $4$ (basis 2).

\textbf{AI-assistance disclosure.}  During the preparation of this
work the author used GitHub Copilot (Microsoft) and Anthropic Claude
in order to assist with code generation, numerical computations, and
manuscript drafting.  After using these tools, the author reviewed
and edited the content as needed and takes full responsibility for
the content of this article.  GitHub Copilot was used as the
execution agent for all enumerations, PSLQ classifications, and
symbolic computations; the agent's outputs are deterministic
functions of the stated scripts.  Anthropic Claude
(\texttt{claude.ai}) was used as a critique and drafting partner
for the manuscript text.  All mathematical claims in this paper
were verified by the author against the underlying numerical and
symbolic data.

\textbf{Author identifier.}  ORCID
\href{https://orcid.org/0009-0000-6192-8273}%
{0009-0000-6192-8273}.

%% ==================================================================
\begin{thebibliography}{99}
%% ==================================================================

\bibitem{TRINote}
Papanokechi.
\emph{The Transcendental Ratio Identity: A Conjecture on
Degree-$(2,1)$ Polynomial Continued Fractions.}
Companion announcement note.
Zenodo preprint,
DOI~\href{https://doi.org/10.5281/zenodo.19801038}%
{10.5281/zenodo.19801038}, 2026.

\bibitem{F25Census}
Papanokechi.
\emph{The F$(2,5)$ Census and the Pigeonhole Falsification.}
SIARC relay bridge session
\texttt{2026-04-29/T2B-F25-FALSIFICATION},
\href{https://github.com/papanokechi/siarc-relay-bridge}%
{papanokechi/siarc-relay-bridge}, commit \texttt{6e28269}, 2026.

\bibitem{SprintW1Report}
Papanokechi.
\emph{Sprint W1 --- Normal Form: Indicial Closed Form and
Compatibility Ideal.}
SIARC relay bridge session
\texttt{2026-04-27/SPRINT-W1-NORMAL-FORM}, 2026.

\bibitem{SprintW2Report}
Papanokechi.
\emph{Sprint W2 --- Trans Indicial Survey: PSLQ Identification of
$|\mu_+/\mu_-|$.}
SIARC relay bridge session
\texttt{2026-04-27/SPRINT-W2-TRANS-INDICIAL-SURVEY}, 2026.

\bibitem{SprintW3Report}
Papanokechi.
\emph{Sprint W3 --- Sub-leading Balance: Resonance Test and
Convergence Rate.}
SIARC relay bridge session
\texttt{2026-04-27/SPRINT-W3-SUBLEADING-BALANCE}, 2026.

\bibitem{ThronWaadeland}
W.~B.~Jones and W.~J.~Thron,
\emph{Continued Fractions: Analytic Theory and Applications},
Encyclopedia of Mathematics and its Applications, vol.~11,
Addison-Wesley, Reading, MA, 1980.

\bibitem{Worpitzky}
J.~Worpitzky,
\emph{Untersuchungen \"uber die Entwickelung der monodromen und
monogenen Functionen durch Kettenbr\"uche},
Programm Friedrichs-Gymnasium und Realschule, Berlin, 1865.

\bibitem{Brouncker}
W.~Brouncker, in J.~Wallis,
\emph{Arithmetica Infinitorum}, Oxford, 1656.

\bibitem{Schwarz}
H.~A.~Schwarz,
\emph{\"Uber diejenigen F\"alle, in welchen die Gauss\-ische
hypergeometrische Reihe eine algebraische Function ihres vierten
Elementes darstellt},
J.\ Reine Angew.\ Math.\ \textbf{75} (1873), 292--335.

\bibitem{RamanujanMachine}
G.~Raayoni et~al.,
\emph{Generating conjectures on fundamental constants with the
Ramanujan Machine},
Nature \textbf{590} (2021), 67--73.

\bibitem{vanderPutSinger}
M.~van der Put and M.~F.~Singer,
\emph{Galois Theory of Difference Equations},
Lecture Notes in Mathematics, vol.~1666, Springer, Berlin, 1997.

\bibitem{JimboSakai}
M.~Jimbo and H.~Sakai,
\emph{A $q$-analog of the sixth Painlevé equation},
Lett.\ Math.\ Phys.\ \textbf{38} (1996), 145--154.

\bibitem{PSLQ}
H.~R.~P.~Ferguson, D.~H.~Bailey, and S.~Arno,
\emph{Analysis of PSLQ, an integer relation finding algorithm},
Math.\ Comp.\ \textbf{68} (1999), 351--369.

\end{thebibliography}

\end{document}
